\documentclass[11pt,a4paper]{article}
\usepackage[T1]{fontenc}
\usepackage{lmodern}
\usepackage[margin=2.5cm]{geometry}
\usepackage{amsmath,amssymb,amsthm}
\usepackage{booktabs}
\usepackage{microtype}
\usepackage{xcolor}
\usepackage{graphicx}
\usepackage[colorlinks=true,linkcolor=blue!50!black,citecolor=blue!50!black,urlcolor=blue!50!black]{hyperref}

\numberwithin{equation}{section}
\newtheorem{proposition}{Proposition}[section]

\newcommand{\dd}{\mathrm{d}}
\newcommand{\upar}{u_\parallel}
\newcommand{\uperp}{u_\perp}
\newcommand{\fh}{\hat f}
\newcommand{\arsinh}{\operatorname{arsinh}}
\newcommand{\arcosh}{\operatorname{arcosh}}
\newcommand{\Kex}{\mathcal K_{\mathrm{ex}}}
\newcommand{\sH}{\mathscr H}
\usepackage{mathrsfs}

\title{Relativistic radiation collisions:\\exact kernel, two-mode closure and radiation drag}
\author{In the notation of Kasilov and Kernbichler (1996)}
\date{28 October 2026}

\begin{document}
\maketitle

\begin{abstract}
The linearised radiation collision operator of Kasilov and Kernbichler~\cite{KK96}
is rederived from a photon-mode variational principle. Summation over cyclotron
harmonics and integration over photon angles give the differential (diffusion)
tensor in closed form. The integral kernel is reduced to an absolutely
convergent special-function series with explicit truncation bounds; its
absorption thresholds produce local $\delta^{3/2}$ cusps. Restricting the
self-consistent photon perturbation to temperature and parallel-drift modes
keeps the exact diffusion tensor and turns the integral part into a rank-two
operator whose two normalisation constants are ratios of modified Bessel
functions. This closure conserves particles, energy and parallel momentum,
dissipates entropy, is the minimal positive correction with these properties,
and differs from the full operator by an exact, computable amount. Separately,
optically thin radiation drag has closed relativistic characteristics and an
explicit propagator. All algebraic steps are checked exactly in the fortsym
computer algebra system, and the main analytic results numerically.
\end{abstract}

\paragraph{Results.}
For the two-mode photon model the full kernel is
\begin{equation}
\begin{gathered}
R_2^{ij'}(y,y')=\frac{D^{i\gamma}(y)\,D^{\gamma j'}(y')}{G_E}
+\frac{D^{i\upar}(y)\,D^{\upar j'}(y')}{G_P},\\
G_E=\frac{4n_e\nu_0}{3\mu}\frac{K_3(\mu)}{K_2(\mu)},\qquad
G_P=\frac{4n_e\nu_0}{15\mu}\frac{K_1(\mu)+K_3(\mu)}{K_2(\mu)},
\end{gathered}
\tag{R}\label{eq:result}
\end{equation}
with $\mu=mc^2/T$, $\nu_0=e^2\omega_{ce}^2/(mc^3\mu)$ and the closed tensor
$D^{ij}$ of Eq.~\eqref{eq:Dclosed}. Applied to a distribution it needs two
velocity moments and no photon integration. For the unchanged sharp-resonance
model the kernel remains a convergent series, Eq.~\eqref{eq:Rkernel}.

\tableofcontents

%======================================================================
\section{Setting and notation}

The operator is the radiation-mediated collision operator of
Ref.~\cite{KK96}, Sec.~IV and Eqs.~(38)--(55), not the Coulomb operator.
Gaussian units are used; $e$ is the magnitude of the electron charge and
temperature is in energy units. The symbol $\mu$ is the inverse
temperature, not the magnetic moment:
\begin{gather}
\mathbf p=mc\,\mathbf u,\quad u=|\mathbf u|,\quad \gamma=\sqrt{1+u^2},\quad
\mu=\frac{mc^2}{T},\quad
\upar=u\cos\chi,\quad \uperp=u\sin\chi,\\
f_0=\frac{n_e\mu}{4\pi(mc)^3K_2(\mu)}e^{-\mu\gamma},\qquad \fh=\frac{\delta f}{f_0}.
\end{gather}
$\chi$ is the particle pitch angle, $\theta$ the photon propagation angle to
$\mathbf B$, and $K_\nu$ the modified Bessel function of the second kind. The
kernel is written in the coordinates
\begin{equation}
y=(\gamma,\upar),\qquad \uperp^2=\gamma^2-\upar^2-1,\qquad
\dd^3p=2\pi(mc)^3\gamma\,\dd\gamma\,\dd\upar=\mathcal J\,\dd\gamma\,\dd\upar .
\end{equation}
Ref.~\cite{KK96} sometimes writes $J=(mc)^3\gamma$, without the $2\pi$ from
the gyrophase; the factor cancels in the divergence but not in momentum
integrals, and is kept here. Electron and photon phase-space integrals and the
Rayleigh--Jeans intensity are
\begin{equation}
\int\dd P=\sum_{n=-\infty}^{\infty}\int\dd^3p,\qquad
\int\dd K=\sum_{M=O,X}\int_0^\infty\dd\omega\int_{4\pi}\dd\Omega,\qquad
I_{\mathrm{RJ}}=\frac{\omega^2T}{8\pi^3c^2},\qquad C_0=\frac{4\pi^2e^2}{m^2c^3}.
\end{equation}
A photon mode is $(M,\omega,\theta,\phi)$; for $\omega>0$ only harmonics
$n\ge1$ resonate. All harmonics feeding one photon mode enter its
normalisation together.

\emph{Assumptions.} Locally homogeneous field, gyrotropic electrons, classical
soft-photon emission and absorption with vacuum O/X polarisation, small
departures from a Maxwell--J\"uttner distribution, and a quasi-static local
photon response. Spatial transport and wall reflection belong to the separate
radiation-transfer problem~\cite{KK96b}. Rayleigh--Jeans applies to the
interacting (low-energy) part of the spectrum.

\emph{Related work.} Radiative electron relaxation in strong fields was treated
by Akhiezer, Peletminskii and coworkers~\cite{Akh62,Akh63,Pel63}; their fast
photon quasi-equilibrium may depend on the non-equilibrium electrons.
Analytic thermal absorption~\cite{Bertelli05}, guiding-centre radiation
reaction~\cite{Hirvijoki15}, optically thin cyclotron cooling~\cite{Kennedy21},
self-absorption-driven thermalisation~\cite{Ghisellini88} and mixed-opacity
models~\cite{Ochs24a,Ochs24b} address related regimes. Conservative
finite-rank field-particle corrections are established for Coulomb
collisions~\cite{Abel08,Sugama09}; the projection used here is the analogue
in photon space.

%======================================================================
\section{Variational form}\label{sec:variational}

With
\begin{equation}
s=\frac{\omega_{ce}}{\omega},\qquad \eta=\cos\theta,\qquad \sigma=1-\eta^2,\qquad
r_n=\gamma-ns-\eta\upar,
\end{equation}
the derivative and polarisation factors of Ref.~\cite{KK96} are
\begin{gather}
\hat L_n=\frac{ns}{\uperp}\partial_{\uperp}+\eta\,\partial_{\upar},\qquad
D_n^{(X)}=\frac{\uperp^2}{\gamma^2}J_n'(\lambda)^2,\qquad
D_n^{(O)}=\frac{(\gamma\eta-\upar)^2}{\gamma^2\sigma}J_n(\lambda)^2,\\
\lambda=\frac{n\uperp\sqrt\sigma}{\gamma-\eta\upar}=\frac{\uperp\sqrt\sigma}{s}
\quad\text{on resonance.}
\end{gather}
On resonance ($r_n=0$) the derivative becomes harmonic-independent in
$(\gamma,\upar)$:
\begin{equation}
\hat L_n=\partial_\gamma+\eta\,\partial_{\upar}\equiv L_\eta,\qquad L_\eta\gamma=1,\qquad L_\eta\upar=\eta .
\label{eq:Leta}
\end{equation}
Define the resonant weight and its moments
\begin{equation}
A_n^{(M)}=\frac{\gamma}{\omega}D_n^{(M)}\delta(r_n)f_0I_{\mathrm{RJ}},\qquad
N=\int\dd P\,A,\qquad B_{\fh}=\int\dd P\,A\,L_\eta\fh .
\end{equation}
With the photon response $q=\mu\,\delta I/I_{\mathrm{RJ}}$ the linearised
quasilinear flux is proportional to $A(L_\eta\fh-q)$, and the quasi-static photon
balance at each mode gives $q_*=B_{\fh}/N$. Equivalently, $q_*$ minimises
\begin{equation}
\mathcal Q[\fh,q]=C_0\int\dd K\int\dd P\,A\,(L_\eta\fh-q)^2 ,
\label{eq:Q}
\end{equation}
and completing the square mode by mode gives the dissipation
\begin{equation}
\Kex[\fh]=\min_q\mathcal Q[\fh,q]
=C_0\int\dd K\left\{\int\dd P\,A\,(L_\eta\fh)^2-\frac{B_{\fh}^2}{N}\right\}\ge0 .
\label{eq:Kex}
\end{equation}
The first term is the differential part, the second the integral correction;
bilinear forms follow by polarisation, and
$\Kex[\hat g,\fh]=-\int\dd^3p\,\hat g\,\hat L_{\mathrm{EM}}(f_0\fh)$ with the
operator of Eq.~\eqref{eq:operator}.

\paragraph{Conservation and null space.}
For $\psi\in\{1,\gamma,\upar\}$, $L_\eta\psi$ is constant at every mode, so
$\Kex[\psi,\fh]=0$: with vanishing boundary flux, number, energy and parallel
momentum are conserved. Conversely, let $\Kex[\fh]=0$ for smooth $\fh$. Then
$L_\eta\fh$ is constant along each resonance line $\gamma-\eta\upar=ns$, whose
tangent is $\eta\partial_\gamma+\partial_{\upar}$:
\begin{equation}
\eta\fh_{\gamma\gamma}+(1+\eta^2)\fh_{\gamma\upar}+\eta\fh_{\upar\upar}=0 .
\end{equation}
Holding for an interval of $\eta$, this forces $\fh_{\gamma\upar}=0$ and
$\fh_{\gamma\gamma}+\fh_{\upar\upar}=0$, so locally
$\fh=c_0+c_E\gamma+c_P\upar+\tfrac{c_2}{2}(\gamma^2-\upar^2)$. The last term has
$L_\eta\fh=c_2ns$ besides the affine part; two harmonics resonating with one
photon mode require $c_2=0$. Hence on the connected physical interior
\begin{equation}
\ker\Kex=\operatorname{span}\{1,\gamma,\upar\}.
\end{equation}
Keeping only one harmonic adds the spurious invariant $\gamma^2-\upar^2$, so
truncation to the fundamental is not a uniformly valid substitute.

%======================================================================
\section{The diffusion tensor}\label{sec:D}

Eliminating the frequency with the resonance delta function
($\int\dd\omega\,\delta(r_n)=\omega_n^2/(n\omega_{ce})$,
$\omega_n=n\omega_{ce}/(\gamma-\eta\upar)$) and integrating over the photon
azimuth gives, with
\begin{equation}
\nu_0=\frac{e^2\omega_{ce}^2}{mc^3\mu},\qquad a=\frac{2\nu_0}{3},\qquad v=(1,\eta)^{\mathsf T},
\end{equation}
\begin{equation}
D^{ij}=\nu_0\int_{-1}^{1}\dd\eta\,\frac{\gamma}{(\gamma-\eta\upar)^3}\sum_{n\ge1}n^2\bigl(D_n^{(X)}+D_n^{(O)}\bigr)v^iv^j .
\label{eq:Dsum}
\end{equation}
Each integrand is positive semidefinite; for $\uperp>0$ the continuum of
photon directions makes $D$ positive definite.

\paragraph{Harmonic sums.} For $|z|<1$,
\begin{equation}
\sum_{n\ge1}n^2J_n(nz)^2=\frac{z^2(4+z^2)}{16(1-z^2)^{7/2}},\quad
\sum_{n\ge1}n^2J_n'(nz)^2=\frac{4+3z^2}{16(1-z^2)^{5/2}} .
\label{eq:sums}
\end{equation}
\emph{Proof.} Under $\vartheta=\varphi-z\sin\varphi$, the Kapteyn expansions give
\[
\frac{1}{1-z\cos\varphi}=1+2\sum_{n\ge1}J_n(nz)\cos n\vartheta,\qquad
\frac{\sin\varphi}{1-z\cos\varphi}=2\sum_{n\ge1}J_n'(nz)\sin n\vartheta .
\]
Parseval's
identity applied to their $\vartheta$-derivatives, transformed back to
$\varphi$, gives
\begin{equation}
\sum_{n\ge1}n^2J_n(nz)^2=\frac1{4\pi}\int_0^{2\pi}\frac{z^2\sin^2\varphi}{(1-z\cos\varphi)^5}\dd\varphi,\qquad
\sum_{n\ge1}n^2J_n'(nz)^2=\frac1{4\pi}\int_0^{2\pi}\frac{(\cos\varphi-z)^2}{(1-z\cos\varphi)^5}\dd\varphi .
\end{equation}
Writing the numerators as polynomials in $\Delta=1-z\cos\varphi$ and using
$\int_0^{2\pi}(\lambda-z\cos\varphi)^{-k}\dd\varphi=\frac{(-1)^{k-1}}{(k-1)!}\partial_\lambda^{k-1}\frac{2\pi}{\sqrt{\lambda^2-z^2}}$
at $\lambda=1$ gives Eq.~\eqref{eq:sums}. \qed

In Eq.~\eqref{eq:Dsum}, $z=\uperp\sqrt\sigma/(\gamma-\eta\upar)<1$. With
$q=\gamma-\eta\upar$ and
\begin{equation}
Q=q^2-\uperp^2(1-\eta^2)=u^2\Bigl(\eta-\frac{\gamma\upar}{u^2}\Bigr)^2+\frac{\uperp^2}{u^2}>0,
\end{equation}
so that $1-z^2=Q/q^2$, the harmonic sum leaves
\begin{equation}
D^{ij}=\frac{\nu_0\uperp^2}{16\gamma}\int_{-1}^{1}\dd\eta\,v^iv^j
\left[\frac{4q^2+3\uperp^2(1-\eta^2)}{Q^{5/2}}+\frac{(\gamma\eta-\upar)^2\bigl(4q^2+\uperp^2(1-\eta^2)\bigr)}{Q^{7/2}}\right].
\label{eq:Deta}
\end{equation}

\paragraph{Angular integration.} The integrands are polynomials in $\eta$
over $Q^{5/2}$ and $Q^{7/2}$, with $Q$ quadratic in $\eta$. Each has an
antiderivative $R(\eta)Q^{-(k-2)/2}+c\,L(\eta)$ with polynomial $R$ and
\begin{equation}
L(\eta)=\frac1u\arsinh\frac{u^2\eta-\gamma\upar}{\uperp},\qquad L'=Q^{-1/2},\qquad L(1)-L(-1)=\frac{2\arsinh u}{u} .
\end{equation}
At the end points $\sqrt{Q(\pm1)}=\gamma\mp\upar$. The logarithmic terms cancel
for $j=0,1$ and survive for $j=2$ (Appendix~\ref{app:angular}). With
\begin{align}
F(u)&=10+\frac{3}{u^2}-\Bigl(\frac{9}{\gamma u}+\frac{3\gamma}{u^3}\Bigr)\arsinh u,\\
G(u)&=10+\frac{15}{u^2}-\Bigl(\frac{21}{\gamma u}+\frac{15\gamma}{u^3}\Bigr)\arsinh u,\qquad
H(u)=u^2-1-\frac{G(u)}{16},
\end{align}
the result is
\begin{equation}
\boxed{\;D^{\gamma\gamma}=a\uperp^2,\qquad
D^{\gamma\upar}=D^{\upar\gamma}=a\frac{\upar\uperp^2}{\gamma},\qquad
D^{\upar\upar}=a\frac{\uperp^2}{u^2}\Bigl[\frac{F(u)}{16}+H(u)\frac{\upar^2}{u^2}\Bigr].\;}
\label{eq:Dclosed}
\end{equation}

\paragraph{Original coordinates.} In $(p,\chi)$,
\begin{gather}
D^{pp}=\frac{2e^2\omega_{ce}^2T}{3c^3}\gamma^2\sin^2\chi,\qquad
D^{p\chi}=D^{\chi p}=\frac{2e^2\omega_{ce}^2T}{3c^3p}\sin\chi\cos\chi,\\
D^{\chi\chi}=\frac{e^2\omega_{ce}^2T}{24c^3p^2}\bigl[F(u)-G(u)\cos^2\chi\bigr],
\end{gather}
which are Eqs.~(48)--(50) of Ref.~\cite{KK96}; Eq.~\eqref{eq:Dclosed} is their
contravariant transformation to $(\gamma,\upar)$. The apparent poles at $u=0$
cancel:
\begin{equation}
F=\tfrac{32}{5}u^2-\tfrac{176}{35}u^4+\tfrac{64}{15}u^6-\tfrac{4352}{1155}u^8+O(u^{10}),\qquad
H=\tfrac{27}{35}u^4-\tfrac{68}{105}u^6+\tfrac{656}{1155}u^8+O(u^{10}),
\label{eq:FHseries}
\end{equation}
which is the stable way to evaluate them near the origin.

\begin{figure}[tb]
\centering
\includegraphics[width=\textwidth]{fig_tensor.pdf}
\caption{(a) The coefficients $F/u^2$ and $H/u^4$ evaluated with the series \eqref{eq:FHseries} below $u=0.1$ (thick); the thin curves are the direct closed forms in double precision, which lose all digits for $u\lesssim10^{-2}$ by cancellation. Dotted: the limits $32/5$ and $27/35$. (b) Parallel diffusion $D^{\upar\upar}/(au^2)$ against $\xi=\cos\chi$ at fixed $u$, compared with $D^{\gamma\gamma}/(au^2)=1-\xi^2$. Relativistic speeds move parallel diffusion from perpendicular towards oblique pitch angles.}
\label{fig:tensor}
\end{figure}

%======================================================================
\section{The integral kernel}\label{sec:kernel}

The operator in divergence form is
\begin{equation}
\hat L_{\mathrm{EM}}\delta f=\frac1{\mathcal J}\partial_i\Bigl\{\mathcal Jf_0\Bigl[D^{ij}\partial_j\fh-\int\dd^3p'\,f_0'R^{ij'}(y,y')\partial_{j'}\fh'\Bigr]\Bigr\},
\label{eq:operator}
\end{equation}
where a primed index refers to the coordinate at the primed momentum. Two
resonances $\gamma-\eta\upar=ns$ and $\gamma'-\eta\upar'=n's$ share the photon
\begin{equation}
\eta_*=\frac{n'\gamma-n\gamma'}{\Delta},\qquad s_*=\frac{\upar\gamma'-\gamma\upar'}{\Delta},\qquad \Delta=n'\upar-n\upar',
\end{equation}
and $|\partial(r_n,r'_{n'})/\partial(s,\eta)|=|\Delta|$. Only $|\eta_*|<1$, $s_*>0$
contribute. The photon measure is $C_0\,2\pi\,(\omega_{ce}\dd s/s^2)\,I_{\mathrm{RJ}}/\omega=\nu_0\,\dd s\,\dd\eta/s^3$,
so for non-degenerate pairs
\begin{equation}
\boxed{\;R^{ij'}(y,y')=P_0\sum_{M=O,X}\;\sum_{n,n'\ge1}^{\mathrm{adm}}
\frac{\gamma\gamma'D_n^{(M)}D_{n'}^{\prime(M)}}{s_*^3|\Delta|\,S_M(s_*,\eta_*)}\,v^i(\eta_*)v^{j'}(\eta_*),\qquad
P_0=\frac{2e^2\omega_{ce}^2K_2(\mu)}{mc^3n_e\mu^2},\;}
\label{eq:Rkernel}
\end{equation}
equivalent to Eq.~(51) of Ref.~\cite{KK96}. The factor $s^{-3}$ collects the
frequency measure and $I_{\mathrm{RJ}}/\omega$.

\paragraph{Absorption normalisation.} $S_M$ is defined by
$N=(I_{\mathrm{RJ}}/\omega)(n_e\mu/2K_2)\,S_M$, i.e.\
$S_M=\sum_k\int\dd\gamma\,\dd\upar\,\gamma^2D_k^{(M)}\delta(\gamma-ks-\eta\upar)e^{-\mu\gamma}$.
A harmonic $k$ absorbs only if $(ks)^2>\sigma$; its resonance line is
\begin{equation}
h_k=\sqrt{(ks)^2-\sigma},\quad \gamma=\frac{ks+\eta h_kx}{\sigma},\quad
\upar=\frac{ks\eta+h_kx}{\sigma},\quad \uperp=\frac{h_k}{\sqrt\sigma}\sqrt{1-x^2},\quad -1<x<1,
\end{equation}
on which $\gamma\eta-\upar=-h_kx$, $\lambda=b_k\sqrt{1-x^2}$ and
$\dd\upar=h_k\dd x/\sigma$. Hence
\begin{gather}
S_M(s,\eta)=\sum_{k:(ks)^2>\sigma}\frac{h_k^3}{\sigma^2}e^{-\mu ks/\sigma}F_k^{(M)}(a_k,b_k),\qquad
a_k=\frac{\mu\eta h_k}{\sigma},\quad b_k=\frac{h_k}{s},\label{eq:SM}\\
F_k^{(O)}(a,b)=\int_{-1}^{1}e^{-ax}x^2J_k\bigl(b\sqrt{1-x^2}\bigr)^2\dd x,\qquad
F_k^{(X)}(a,b)=\int_{-1}^{1}e^{-ax}(1-x^2)J_k'\bigl(b\sqrt{1-x^2}\bigr)^2\dd x .
\end{gather}
The $x$ integrals are entire in $a,b$ and have the series (Appendix~\ref{app:series})
\begin{gather}
F_k^{(O)}=\sum_{\ell\ge0}(-1)^\ell A_{k\ell}\Bigl(\frac b2\Bigr)^{2k+2\ell}T_{k+\ell}(a),\qquad
F_k^{(X)}=\frac1{b^2}\sum_{\ell\ge0}(-1)^\ell B_{k\ell}\Bigl(\frac b2\Bigr)^{2k+2\ell}M_{k+\ell}(a),\label{eq:Fseries}\\
A_{k\ell}=\frac{(2k+2\ell)!}{\ell!\,(2k+\ell)!\,(k+\ell)!^2},\qquad
B_{k\ell}=\bigl[2(k+\ell)^2-k^2\bigr]A_{k\ell}-4A_{k,\ell-1},\quad A_{k,-1}=0,\\
M_j(a)=\int_{-1}^1e^{-ax}(1-x^2)^j\dd x=B\bigl(\tfrac12,j+1\bigr)\,{}_0F_1\bigl(;j+\tfrac32;\tfrac{a^2}4\bigr),\\
T_j(a)=\int_{-1}^1x^2e^{-ax}(1-x^2)^j\dd x=B\bigl(\tfrac32,j+1\bigr)\,{}_1F_2\bigl(\tfrac32;\tfrac12,j+\tfrac52;\tfrac{a^2}4\bigr).
\end{gather}
These converge absolutely for all finite arguments; they remove quadrature from
the coefficients, but the harmonic sums remain infinite, and cancellation at
large arguments needs extended precision.

\paragraph{Harmonic tail.} For $s>0$, $|\eta|<1$: $|J_k|,|J_k'|\le1$, $h_k\le ks$ and
$ks-|\eta|h_k\ge ks(1-|\eta|)$ give, with $q=e^{-\mu s/(1+|\eta|)}<1$,
\begin{equation}
0\le S_X-S_{X,K}\le\frac{4s^3}{3\sigma^2}\sum_{k>K}k^3q^k,\qquad
0\le S_O-S_{O,K}\le\frac{2s^3}{3\sigma^2}\sum_{k>K}k^3q^k .
\end{equation}
$S_M$ is strictly positive at every physical mode.

\paragraph{Limits.} On the propagation axis only the fundamental survives, with
$D_1^{(X)},D_1^{(O)}\to\uperp^2/(4\gamma^2)$, $\gamma-\upar=s$ and
$\uperp^2=2s(\gamma-\gamma_{\min})$, $\gamma_{\min}=(s+s^{-1})/2$:
\begin{equation}
S_X(s,\pm1)=S_O(s,\pm1)=\frac{s}{2\mu^2}\exp\Bigl[-\frac\mu2(s+s^{-1})\Bigr].
\label{eq:axis}
\end{equation}
In the dipole limit ($\tau=1/\mu\ll1$, fundamental only, Doppler resonance
$1-s-\eta\upar=0$, away from $\eta=0$), $D_1^{(X)}\simeq\uperp^2/4$,
$D_1^{(O)}\simeq\eta^2\uperp^2/4$ and the perpendicular average at fixed $\upar$,
\begin{equation}
\Pi F=\frac1{2\tau^2}\int_0^\infty w^3e^{-w^2/(2\tau)}F(w,\upar)\,\dd w,\qquad \Pi[w^{2j}]=(2\tau)^j(j+1)!,
\end{equation}
makes all polynomial moments explicit. This one-harmonic limit retains the
spurious invariant of Sec.~\ref{sec:variational}.

%======================================================================
\section{Kernel properties}\label{sec:properties}

\paragraph{Resonance degeneracy.} For fixed $y,n,n'$ let $\rho=n'/n$,
$a_1=\gamma'-\rho\gamma$, $b_1=\upar'-\rho\upar$. Eliminating the frequency
leaves $\delta(a_1-\eta b_1)$, so near a coincident resonance a harmonic pair
contributes $\mathbf 1_{|a_1|<|b_1|}H(a_1/b_1,b_1)/|b_1|$ with $H$ bounded.
The apparent pole is integrable,
$\int_{-\epsilon}^{\epsilon}\dd b_1\int_{-|b_1|}^{|b_1|}\dd a_1/|b_1|=4\epsilon$,
and the critical set $b_1=0$ carries no point mass; the kernel is a density
defined almost everywhere.

\paragraph{Row sum.} Integrating the primed numerator of Eq.~\eqref{eq:Rkernel}
against $f_0'\dd^3p'$ reproduces $S_M$, so
\begin{equation}
\int\dd^3p'\,f_0'R^{ij'}(y,y')=D^{ij}(y).
\label{eq:rowsum}
\end{equation}
Each harmonic pair is a non-negative multiple of $vv^{\mathsf T}$, hence
$0\le R^{\upar\upar'}\le R^{\gamma\gamma'}$, $|R^{\gamma\upar'}|\le R^{\gamma\gamma'}$, and
$\int\dd^3p'f_0'R^{\gamma\gamma'}=a\uperp^2<\infty$. By Tonelli's theorem the full
positive harmonic series converges absolutely almost everywhere and in
weighted $L^1_{\mathrm{loc}}$. With exchange symmetry, Eq.~\eqref{eq:rowsum} gives the
weak form
\begin{equation}
\Kex[\hat g,\fh]=\frac12\iint\dd^3p\,\dd^3p'\,f_0f_0'\,
[\nabla\hat g(y)-\nabla\hat g(y')]^{\mathsf T}R(y,y')[\nabla\fh(y)-\nabla\fh(y')],
\end{equation}
with the same $(\gamma,\upar)$ components on both factors.

\paragraph{Harmonic truncation.} Let $w=\uperp>0$ and
\begin{equation}
\alpha(w)=\arsinh\frac1w-\frac1{\sqrt{1+w^2}},\qquad q_w=e^{-2\alpha(w)} .
\end{equation}
Over photon directions $z\le z_m=w/\sqrt{1+w^2}$, because
$q^2-(1+w^2)\sigma=(\gamma\eta-\upar)^2\ge0$. Shifting the contour of
$J_n(x)=\frac1{2\pi}\int_{-\pi}^{\pi}e^{i(x\sin\vartheta-n\vartheta)}\dd\vartheta$ by
$-i\beta_m$, $\cosh\beta_m=1/z_m$, and using
$\frac nxJ_n=\frac1{2\pi}\int\cos\vartheta\,e^{i(x\sin\vartheta-n\vartheta)}\dd\vartheta$ gives
for $z\le z_m$
\begin{equation}
|J_n(nz)|\le\frac{z}{z_m}e^{-n\alpha(w)},\qquad |J_n'(nz)|\le\frac{1}{z_m}e^{-n\alpha(w)} .
\end{equation}
With $(\gamma\eta-\upar)^2\le q^2$ this bounds
$n^2(D_n^{(X)}+D_n^{(O)})\le2(1+w^2)n^2q_w^n/\gamma^2$, and
$\int_{-1}^1\dd\eta\,q^{-3}=2\gamma/(1+w^2)^2$ gives
\begin{equation}
D^{\gamma\gamma}_{>N}(y)\le\frac{4\nu_0}{1+w^2}\sum_{n>N}n^2q_w^n,\qquad
\sum_{n\ge m}n^2q^n=q^m\Bigl[\frac{m^2}{1-q}+\frac{2mq}{(1-q)^2}+\frac{q(1+q)}{(1-q)^3}\Bigr].
\label{eq:tail}
\end{equation}
The thermal tail is exact: with $A=\sqrt{1+W^2}$ and
$\int\dd\upar\,e^{-\mu\sqrt{A^2+\upar^2}}=2AK_1(\mu A)$,
\begin{equation}
T_W=\int_{w>W}\dd^3p\,f_0w^2=\frac{n_e}{K_2(\mu)}\Bigl[W^2A^2K_2(\mu A)+\frac{2A^3}{\mu}K_3(\mu A)\Bigr].
\end{equation}
Since $q_w$ increases with $w$, splitting at $w=W$ gives
$\int\dd^3p\,f_0D^{\gamma\gamma}_{>N}\le4n_e\nu_0\sum_{n>N}n^2q_W^n+aT_W$, which
tends to zero, e.g.\ for $W\propto N^{1/4}$ ($\alpha(W)\simeq1/(3W^3)$). If
$|\partial_\gamma\fh|+|\partial_{\upar}\fh|\le M$, truncating every mode
consistently at harmonic $N$ satisfies
\begin{equation}
0\le\Kex[\fh]-\mathcal K_N[\fh]\le4M^2\int\dd^3p\,f_0D^{\gamma\gamma}_{>N}.
\end{equation}
At each mode $\mathcal K=C_0\,N\operatorname{Var}_A(L_\eta\fh)$; adding weight
$\Delta N$ cannot decrease the sum of squared deviations, and evaluating it at
the old mean, with $|L_\eta\fh-\bar L|\le2M$, increases it by at most
$4M^2\Delta N$. These bounds concern exact arithmetic, not floating-point
cancellation.

\begin{figure}[tb]
\centering
\includegraphics[width=0.62\textwidth]{fig_harmonics.pdf}
\caption{Relative harmonic tail $D^{\gamma\gamma}_{>N}/D^{\gamma\gamma}$ of Eq.~\eqref{eq:Dsum} at $\upar=0.3$ (Gauss--Legendre in $\eta$, exact total from Eq.~\eqref{eq:Dclosed}) and the bound \eqref{eq:tail} (dotted, same colour). Convergence is geometric with rate $q_w=e^{-2\alpha(\uperp)}$, which approaches one as $\uperp$ grows; the floor near $10^{-14}$ is double-precision roundoff.}
\label{fig:harmonics}
\end{figure}

%======================================================================
\section{Absorption thresholds}\label{sec:thresholds}

\paragraph{Local cusp.} Let $\delta=s^2-\sigma$. At the opening of the
fundamental absorbing harmonic, Eq.~\eqref{eq:SM} gives
\begin{equation}
S_{X,1}=\frac{\delta^{3/2}}{\sigma^2}e^{-\mu s/\sigma}\int_{-1}^{1}e^{-\mu\eta\sqrt\delta x/\sigma}(1-x^2)
J_1'\Bigl(\frac{\sqrt\delta}{s}\sqrt{1-x^2}\Bigr)^2\dd x
\end{equation}
for $\delta>0$ and zero otherwise. With $J_1(z)=z/2+O(z^3)$ and
$J_1'(z)=1/2+O(z^2)$, odd terms integrate to zero and
\begin{equation}
S_{X,1}=a_X\delta_+^{3/2},\quad a_X\big|_{\delta=0}=\frac{e^{-\mu s/\sigma}}{3\sigma^2},\qquad
S_{O,1}=a_O\delta_+^{5/2},\quad a_O\big|_{\delta=0}=\frac{e^{-\mu s/\sigma}}{15s^2\sigma^2},
\end{equation}
with $a_X,a_O$ analytic in the transverse coordinate. The higher harmonics are
already open, so $B_X=\sum_{k\ge2}S_{X,k}$ is positive and smooth and
\begin{equation}
\frac1{S_X}=\frac1{B_X}-\frac{a_X}{B_X^2}\delta_+^{3/2}+\frac{a_X^2\delta_+^3}{B_X^2(B_X+a_X\delta_+^{3/2})},
\end{equation}
the last term being $C^2$. A harmonic pair with non-zero numerator therefore
contributes a negative $\delta^{3/2}$ cusp to $R^{\gamma\gamma'}$, which the
ordinary mode cannot cancel.

\begin{figure}[tb]
\centering
\includegraphics[width=\textwidth]{fig_absorption.pdf}
\caption{(a) Absorption normalisation $S_M(s,\eta)$, Eq.~\eqref{eq:SM}, at $\eta=0.5$, $\mu=10$. Each harmonic $k$ opens at $ks=\sqrt\sigma$ (dotted) with a kink in $S_X$ and a weaker one in $S_O$. (b) Opening of the fundamental: $S_{X,1}/a_X$ and $S_{O,1}/a_O$ against $\delta=s^2-\sigma$ follow $\delta^{3/2}$ and $\delta^{5/2}$ (dotted), Eqs.~(6.2). The $\delta^{3/2}$ onset of $S_X$ is the origin of the kernel cusp.}
\label{fig:absorption}
\end{figure}

\paragraph{Dense thresholds.} On the counter-streaming slice $y=(\gamma,v)$,
$y'=(\gamma,-v)$, $v>0$, with $N=n+n'$ and $d=n'-n$, the common photon is
$s=2\gamma/N$, $\eta=d\gamma/(Nv)$, and the fundamental threshold lies at
\begin{equation}
\gamma_{N,d}=\frac{Nv}{\sqrt{d^2+4v^2}},\qquad |d|<N,\quad N\equiv d\pmod2 .
\end{equation}
Choosing $d_j/N_j\to v/\gamma_0$ with $N_j\to\infty$ gives $\gamma_{N_j,d_j}\to\gamma_0$,
so these thresholds are dense in $\gamma>\sqrt{1+v^2}$, each at an interior
photon direction $|\eta|=|d|/\sqrt{d^2+4v^2}<1$.

\begin{proposition}[conditional]
Let $U$ be an open momentum region in which the thresholds $\delta_{nn'}=0$
are dense. Suppose that on $U$
$R^{\gamma\gamma'}=\sum_{n,n'}\bigl[G_{nn'}-A_{nn'}\delta_{nn',+}^{3/2}\bigr]$ with
$G_{nn'},A_{nn'}\in C^2(U)$, $A_{nn'}\ge0$, $A_{nn'}>0$ on $\{\delta_{nn'}=0\}$, and
\[
\|G_{nn'}\|_{C^2}+\|A_{nn'}\|_{C^2}\le CN^pK^NW^{2N-8},\qquad KW^2<1,
\]
for $W/2<\uperp,\uperp'<W$ and $N=n+n'$. Then $R^{\gamma\gamma'}$ has no $C^2$
representative on any open subset of $U$; in particular it is not a finite
expression in locally analytic functions with finitely many branch loci.
\end{proposition}
\begin{proof}
Transverse to a threshold, the singular part of the second derivative,
$-\tfrac34A_{nn'}|\nabla\delta_{nn'}|^2\delta_{nn'}^{-1/2}$ on $\delta_{nn'}>0$, is
unbounded wherever $A_{nn'}>0$. The summability bound keeps $\sum G_{nn'}$ in
$C^2$ and every other contribution locally bounded, the singular terms share
one sign and cannot cancel, and every open subset of $U$ meets a threshold.
\end{proof}
The cusp and the threshold locations are established. The derivative estimate
assumed in the proposition requires uniform control of reciprocal absorption
coefficients near the axis and near thresholds and is not proved here; the
row-sum identity and pointwise convergence do not imply it. Whether a kernel
has a ``closed form'' depends on the admissible function class, which must be
stated. Nothing below depends on the proposition.

%======================================================================
\section{Modified physics}\label{sec:modified}

\paragraph{Finite line width.} Convolving the local cusp with
$g_\ell(x)=e^{-x^2/(2\ell^2)}/(\sqrt{2\pi}\ell)$, substituting $t=\ell y$ and
using $\int_0^\infty t^{\nu-1}e^{-t^2/2-zt}\dd t=\Gamma(\nu)e^{z^2/4}D_{-\nu}(z)$~\cite{DLMF},
\begin{equation}
(g_\ell*x_+^{3/2})(x)=\frac{\Gamma(5/2)\ell^{3/2}}{\sqrt{2\pi}}e^{-x^2/(4\ell^2)}D_{-5/2}(-x/\ell),
\end{equation}
where $D_{-5/2}$ is the parabolic-cylinder function. This is analytic in $x$,
so the threshold argument does not carry over; but the broadened kernel still
contains the reciprocal of a mode-dependent absorption sum, which smoothing
does not evaluate.

\paragraph{Broadening must respect detailed balance.} Off resonance,
\begin{equation}
\hat L_n\gamma=\frac{ns+\eta\upar}{\gamma}=1-\frac{r_n}{\gamma},\qquad
\hat L_nf_0+\mu f_0=\mu\frac{r_n}{\gamma}f_0 .
\end{equation}
The right-hand side vanishes only on the sharp resonance. Replacing
$\delta(r_n)$ by a Gaussian or Lorentzian while keeping the drift and response
loses the detailed-balance derivation; the broadening mechanism and its
momentum and energy exchange must be specified. A positive broadened model can
instead be built with $L_\eta$ throughout, keeping the invariants; its
diffusion tensor must then be recomputed.

\paragraph{Other changes.} Planck radiation with recoil changes both the
emission--absorption balance and the normaliser. Refraction adds
frequency-dependent polarisation and dispersion. A synchrotron continuum
replaces the harmonic sums in the ultrarelativistic regime but leaves a
spectrum-dependent absorption problem~\cite{Ochs24a,Ochs24b}. An imposed
thermal photon bath keeps only the differential operator and transfers energy
and momentum to the reservoir. None of these yields a finite kernel. The
useful reduction acts on the photon response rather than on electron
harmonics.

%======================================================================
\section{Two-mode closure}\label{sec:closure}

\paragraph{Photon model.} Restrict the photon response to
\begin{equation}
V=\operatorname{span}\{1,\eta\},\qquad q_2=a_E+a_P\eta .
\end{equation}
In the Rayleigh--Jeans regime this is a radiation field at a slightly
shifted temperature with a small parallel drift,
$\delta I/I_{\mathrm{RJ}}=\delta T/T+(V_\parallel/c)\eta$. The same two coefficients
apply to both polarisations and all interacting frequencies: the model assumes
redistribution to a common spectral and angular equilibrium, not only fast
absorption mode by mode. Photon number is not an invariant. In the
absorption-weighted inner product
\begin{equation}
\langle r,s\rangle_N=C_0\int\dd K\,N\,rs
\end{equation}
all norms are finite, and minimising Eq.~\eqref{eq:Q} over $V$ gives the
orthogonal projection $q_2=Pq_*$. With $v_E=1$, $v_P=\eta$ the normal equations are
\begin{equation}
\mathsf G_{\alpha\beta}a_\beta=b_\alpha[\fh],\qquad
\mathsf G_{\alpha\beta}=\langle v_\alpha,v_\beta\rangle_N=\int\dd^3p\,f_0D^{\alpha\beta},\qquad
b_\alpha[\fh]=\langle v_\alpha,q_*\rangle_N=\int\dd^3p\,f_0D^{\alpha j}\partial_j\fh,
\end{equation}
with $E\leftrightarrow\gamma$, $P\leftrightarrow\upar$. As $f_0$ is even in
$\upar$, $\mathsf G=\operatorname{diag}(G_E,G_P)$ and $q_2=b_E/G_E+(b_P/G_P)\eta$.

\paragraph{Normalisation constants.} With $\langle F\rangle_0=n_e^{-1}\int f_0F\,\dd^3p$
and the representation
$\int_1^\infty e^{-\mu\gamma}(\gamma^2-1)^{\nu-1/2}\dd\gamma=2^\nu\Gamma(\nu+\tfrac12)K_\nu(\mu)/(\sqrt\pi\mu^\nu)$~\cite{DLMF},
differentiation in $\mu$ and the recurrences give
\begin{equation}
\langle u^2\rangle_0=\frac{3K_3(\mu)}{\mu K_2(\mu)},\qquad\langle\uperp^2\rangle_0=\frac{2K_3(\mu)}{\mu K_2(\mu)},\qquad
\boxed{G_E=\frac{2an_e}{\mu}\frac{K_3}{K_2}=\frac{4n_e\nu_0}{3\mu}\frac{K_3(\mu)}{K_2(\mu)}.}
\end{equation}
For the parallel constant, average at fixed $u$ over $\xi=\cos\chi$:
\begin{equation}
\frac12\int_{-1}^{1}D^{\upar\upar}\dd\xi=a\Bigl(\frac F{24}+\frac{2H}{15}\Bigr)
=\frac a{15}\Bigl[3+2u^2-\frac{3\arsinh u}{\gamma u}\Bigr].
\end{equation}
With $u\,\dd u/\gamma=\dd\gamma$ and one integration by parts,
\begin{equation}
\Bigl\langle\frac{\arsinh u}{\gamma u}\Bigr\rangle_0=\frac{\mu}{K_2}\int_1^\infty\arcosh\gamma\,e^{-\mu\gamma}\dd\gamma=\frac{K_0(\mu)}{K_2(\mu)} .
\end{equation}
Using $K_2-K_0=2K_1/\mu$,
\begin{gather}
\boxed{G_P=\frac{an_e}{15}\Bigl[3+\frac{6K_3}{\mu K_2}-\frac{3K_0}{K_2}\Bigr]=\frac{4n_e\nu_0}{15\mu}\frac{K_1(\mu)+K_3(\mu)}{K_2(\mu)},}\\
\frac{G_P}{G_E}=\frac15\Bigl(1+\frac{K_1}{K_3}\Bigr)\to
\begin{cases}2/5,&\mu\to\infty,\\1/5,&\mu\to0.\end{cases}
\end{gather}
No low-temperature expansion is involved.

\begin{figure}[tb]
\centering
\includegraphics[width=\textwidth]{fig_constants.pdf}
\caption{(a) Closure constants $\mu G_E/(n_e\nu_0)=\tfrac43K_3/K_2$ and $\mu G_P/(n_e\nu_0)=\tfrac4{15}(K_1+K_3)/K_2$. (b) Their ratio moves from $2/5$ in the nonrelativistic limit to $1/5$ in the ultrarelativistic limit, with the transition around $T\sim mc^2/10$.}
\label{fig:constants}
\end{figure}

\paragraph{Operator.} Inserting $q_2$ into the electron flux gives the kernel
\eqref{eq:result}, $R_2=d_Ed_E'^{\mathsf T}/G_E+d_Pd_P'^{\mathsf T}/G_P$ with
\begin{equation}
d_E(y)=a\begin{pmatrix}\uperp^2\\ \upar\uperp^2/\gamma\end{pmatrix},\qquad
d_P(y)=a\begin{pmatrix}\upar\uperp^2/\gamma\\ (\uperp^2/u^2)[F/16+H\upar^2/u^2]\end{pmatrix},
\end{equation}
and, in the convention of Eq.~\eqref{eq:operator},
\begin{equation}
\boxed{\hat L_2\delta f=\frac1{\mathcal J}\partial_i\Bigl\{\mathcal Jf_0\Bigl[D^{ij}\partial_j\fh-D^{i\gamma}\frac{b_E[\fh]}{G_E}-D^{i\upar}\frac{b_P[\fh]}{G_P}\Bigr]\Bigr\}.}
\end{equation}
$R_2(y,y')=R_2(y',y)^{\mathsf T}$. An off-diagonal block of $R_2$ need not be
positive semidefinite; positivity of the whole operator follows from its
variational construction. The constants are local in space when $B$, $T$ and
$n_e$ vary on scales compatible with the local approximation.

\paragraph{Field-particle form.} Define
$\sH_\alpha=-(\mathcal Jf_0)^{-1}\partial_i(\mathcal Jf_0D^{i\alpha})$. Direct
differentiation gives
\begin{equation}
\sH_E=a\Bigl[\mu\uperp^2-\frac{2+4\uperp^2}{\gamma}\Bigr],\qquad
\sH_P=-a\upar\Bigl[2-\frac{\mu\uperp^2}{\gamma}-\frac{F(u)}{8u^2}+\Bigl(\frac2{u^2}-\frac{4\upar^2}{u^4}\Bigr)H(u)\Bigr],
\end{equation}
with removable limits at the origin from Eq.~\eqref{eq:FHseries}. For regular
gyrotropic distributions with adequate decay the boundary terms vanish (the
$\uperp=0$ boundary carries no diffusive flux), so $b_\alpha[\fh]=\int\dd^3p\,\delta f\,\sH_\alpha$
and the integral part is the explicit rank-two source
\begin{equation}
\boxed{\hat L_{I,2}\delta f=f_0\Bigl[\frac{\sH_E}{G_E}\int\dd^3p'\,\sH_E'\delta f'+\frac{\sH_P}{G_P}\int\dd^3p'\,\sH_P'\delta f'\Bigr],}
\end{equation}
added to the exact differential part $\hat L_D$.

\begin{figure}[tb]
\centering
\includegraphics[width=\textwidth]{fig_sources.pdf}
\caption{Scalar sources of the closure at $\mu=5$ in the $(\upar,\uperp)$ plane; black contours are $f_0/f_0(0)=e^{-1},e^{-3},e^{-6}$. (a) $\sH_E/a$ is even in $\upar$ and changes sign inside the thermal bulk, so $\int f_0\sH_E=0$; (b) $\sH_P/a$ is odd in $\upar$. Together with $f_0$ these two functions are the complete integral part, Eq.~(8.12).}
\label{fig:sources}
\end{figure}

%======================================================================
\section{Properties of the closure}\label{sec:closureprops}

\paragraph{Structure.} The quadratic form is
\begin{equation}
\mathcal K_2[\fh]=\mathcal K_D[\fh]-\frac{b_E[\fh]^2}{G_E}-\frac{b_P[\fh]^2}{G_P}=\min_{q\in V}\mathcal Q[\fh,q]\ge0,
\end{equation}
so $\hat L_2$ is self-adjoint in the linearised entropy pairing. Since
$b[\gamma]=(G_E,0)^{\mathsf T}$ and $b[\upar]=(0,G_P)^{\mathsf T}$,
$\hat L_2(f_0)=\hat L_2(f_0\gamma)=\hat L_2(f_0\upar)=0$, and with the stated boundary
conditions
\begin{equation}
\int\dd^3p\,(1,mc^2\gamma,mc\upar)\,\hat L_2\delta f=0,\qquad
\frac{\dd}{\dd t}\frac12\int\dd^3p\,f_0\fh^2=-\mathcal K_2[\fh]\le0 .
\end{equation}
This is a linearised statement; it implies nothing about positivity of
large-amplitude distributions. There is no extra null mode: $\mathcal K_2=0$
requires $\partial_\gamma\fh+\eta\partial_{\upar}\fh=a_E+a_P\eta$ on every weighted
resonance, the continuum of $\eta$ forces $\nabla\fh=(a_E,a_P)$, and
$\ker\mathcal K_2=\operatorname{span}\{1,\gamma,\upar\}$. Unlike Sec.~\ref{sec:variational},
$q_2$ is independent of $s$, so even a single harmonic would not add
$\gamma^2-\upar^2$. The row sum $\int\dd^3p'f_0'R_2^{ij'}=D^{ij}$ holds exactly.

\paragraph{Minimality.} Let $\mathcal A=-f_0^{-1}\hat L_Df_0$ act on $L^2(f_0\dd^3p)$
with form $\mathcal K_D$, and $\psi_E=\gamma$, $\psi_P=\upar$; then
$\mathcal A\psi_\alpha=\sH_\alpha$, independent (one even, one odd in $\upar$). If a
positive self-adjoint correction $\mathcal C$ enforces $(\mathcal A-\mathcal C)\psi_\alpha=0$,
then $\langle\psi_\alpha,\mathcal C\psi_\beta\rangle=\mathsf G_{\alpha\beta}$ and
$\langle\psi_\alpha,\mathcal C\fh\rangle=b_\alpha[\fh]$, and Cauchy--Schwarz in the
$\mathcal C$-seminorm gives $\langle\fh,\mathcal C\fh\rangle\ge b^{\mathsf T}\mathsf G^{-1}b$.
The rank-two correction attains equality: it is the minimal positive
correction, in form order, that restores the invariants with $D$ fixed. This is
a structural optimality, not a statement about transport accuracy.

\paragraph{Error identity.} Since $q_*$ is the $A$-mean at each mode,
$\mathcal Q[\fh,q]=\Kex[\fh]+\|q-q_*\|_N^2$, and therefore
\begin{equation}
\boxed{\mathcal K_2[\fh]-\Kex[\fh]=\|(\mathrm{Id}-P)q_*\|_N^2\ge0 .}
\end{equation}
The closure is at least as dissipative as the original operator and exact for
an input precisely when its induced photon response lies in $V$. This is an a
posteriori diagnostic: if $\|(\mathrm{Id}-P)q_*\|_N\le\varepsilon\sqrt{\mathcal K_D[\fh]}$
then $0\le\mathcal K_2-\Kex\le\varepsilon^2\mathcal K_D$. No uniform relative error
holds for arbitrary distributions, and conservation is not an accuracy
estimate. Nested photon spaces $V\subset V_3\subset V_4\subset\cdots$ give
$\mathcal K_D\ge\mathcal K_2\ge\mathcal K_3\ge\mathcal K_4\ge\Kex$; their extra Gram
entries can be precomputed without forming the two-electron kernel.

%======================================================================
\section{Physical interpretation}\label{sec:interpretation}

\paragraph{Redistribution model.} Penalising photon components outside $V$,
\begin{equation}
\mathcal Q_\kappa[\fh,q]=\mathcal Q[\fh,q]+\kappa\|(\mathrm{Id}-P)q\|_N^2,\qquad\kappa\ge0,
\end{equation}
defines a dissipative model of additional photon redistribution whose slow
variables are common temperature and parallel drift. It is a specified model,
not derived from a particular scattering cross section. The minimiser is
$q_\kappa=Pq_*+(\mathrm{Id}-P)q_*/(1+\kappa)$, and
\begin{equation}
\boxed{\mathcal K_\kappa=\frac{\Kex+\kappa\mathcal K_2}{1+\kappa},\qquad
0\le\mathcal K_2-\mathcal K_\kappa=\frac{\|(\mathrm{Id}-P)q_*\|_N^2}{1+\kappa}\le\frac{\mathcal K_D}{1+\kappa}.}
\end{equation}
The strong-redistribution limit is the closure, with an $O(\kappa^{-1})$ form
error. With the equilibrium absorption coefficient $\alpha_0(M,\omega,\eta)$ of
Ref.~\cite{KK96}, $C_0N=\alpha_0I_{\mathrm{RJ}}/(mc^2\mu)$, and the stationary
photon equation
\begin{equation}
0=c\alpha_0(q_*-q)-\kappa c\alpha_0(\mathrm{Id}-P)q
\end{equation}
has the same minimiser. $P$ is orthogonal in the $\alpha_0I_{\mathrm{RJ}}$ weight, so
the added term exchanges no energy or parallel momentum in this description;
its rate is $\nu_{\mathrm{red}}=\kappa c\alpha_0$.

\paragraph{Applicability.} The model is closer to reality than independent
photon modes when an actual mechanism redistributes the interacting spectrum
and angles quickly compared with absorption and with the electron evolution,
and conserves the same slow variables. Optical thickness alone does not
provide this: angular scattering need not thermalise frequency, scattering by
electrons exchanges energy and momentum with them, and walls absorb momentum.
In a dilute tokamak the strong-redistribution condition must be checked.

\paragraph{External bath.} Adding $a^{\mathsf T}\Lambda a$ with a symmetric positive
semidefinite $2\times2$ matrix $\Lambda$ keeps the closure finite:
\begin{equation}
a=(\mathsf G+\Lambda)^{-1}b,\qquad R_\Lambda^{ij'}=D^{i\alpha}(y)[\mathsf G+\Lambda]^{-1}_{\alpha\beta}D^{\beta j'}(y').
\end{equation}
Number is conserved; energy and parallel momentum flow to the reservoir at
the rate given by $b-\mathsf G(\mathsf G+\Lambda)^{-1}b=\Lambda(\mathsf G+\Lambda)^{-1}b$.
Large $\Lambda$ clamps a pre-existing thermal bath and leaves the differential
operator. This is not the vacuum escape limit, which removes the induced
diffusion as well.

%======================================================================
\section{Optically thin limit: radiation drag}\label{sec:drag}

If the optical depth is small over the emitted spectrum and incident radiation
is negligible, reabsorption is absent and only the classical relativistic
radiation drag of Ref.~\cite{KK96}, Eq.~(55), remains (see also
Refs.~\cite{Kennedy21,Hirvijoki15}). Without Coulomb collisions, imposed fields
or spatial variation, set
\begin{equation}
\nu_R=\frac{2e^2\omega_{ce}^2}{3mc^3}=\mu a,\qquad w=\uperp,\qquad v=\upar .
\end{equation}
The characteristics
\begin{equation}
\dot w=-\nu_R\frac{w(1+w^2)}{\gamma},\qquad \dot v=-\nu_R\frac{vw^2}{\gamma},\qquad \dot\gamma=-\nu_Rw^2
\end{equation}
conserve the parallel velocity:
\begin{equation}
\beta=\frac v\gamma,\qquad \gamma_\parallel=(1-\beta^2)^{-1/2},\qquad \gamma=\gamma_\parallel\sqrt{1+w^2}.
\end{equation}
With $\zeta=\arsinh(1/w)$, $\dot\zeta=\nu_R/\gamma_\parallel$ and
\begin{equation}
\boxed{\zeta(t)=\zeta_0+\frac{\nu_Rt}{\gamma_\parallel},\quad w=\operatorname{csch}\zeta,\quad
\gamma=\gamma_\parallel\coth\zeta,\quad v=\beta\gamma_\parallel\coth\zeta .}
\end{equation}
An electron loses $mc^2[\gamma(0)-\gamma(t)]$, carried away by the radiation;
$w=0$ is stationary.

\paragraph{Propagator.} In characteristic coordinates
\begin{equation}
\dd^3p=\mathcal J_\zeta\,\dd\zeta\,\dd\beta,\qquad \mathcal J_\zeta=2\pi(mc)^3\gamma_\parallel^3\operatorname{csch}^2\zeta\coth^2\zeta,
\end{equation}
and particle conservation reads
$\partial_t(\mathcal J_\zeta f)+\partial_\zeta(\nu_R\gamma_\parallel^{-1}\mathcal J_\zeta f)=0$.
The speed is independent of $\zeta$, so with $\zeta_0=\zeta-\nu_Rt/\gamma_\parallel$
\begin{equation}
\boxed{f(\zeta,\beta,t)=\frac{\mathcal J_\zeta(\zeta_0,\beta)}{\mathcal J_\zeta(\zeta,\beta)}f_{\mathrm{in}}(\zeta_0,\beta)
=\Bigl[\frac{\sinh\zeta}{\sinh\zeta_0}\Bigr]^2\Bigl[\frac{\coth\zeta_0}{\coth\zeta}\Bigr]^2f_{\mathrm{in}}(\zeta_0,\beta)}
\end{equation}
for $\zeta_0>0$, and $f=0$ for $\zeta_0\le0$ (no particles arrive from infinite
momentum). This solves the gyrotropic drag problem for arbitrary admissible
initial data in a uniform field. It does not conserve electron energy and does
not solve the reabsorbing problem.

\begin{figure}[tb]
\centering
\includegraphics[width=\textwidth]{fig_drag.pdf}
\caption{Optically thin radiation drag, Eq.~(11.4). (a) Trajectories in $(\upar,\uperp)$ from four initial points, with markers every $\nu_Rt=1$; each electron moves on its line of constant $\beta=\upar/\gamma$ towards $\uperp=0$. (b) Kinetic energy $\gamma-1$ relaxes to $\gamma_\parallel-1$; a purely perpendicular electron (dotted) loses all of it.}
\label{fig:drag}
\end{figure}

%======================================================================
\section{Verification}\label{sec:verification}

\paragraph{Exact algebra.} Every algebraic step is checked with the native
fortsym computer algebra system~\cite{fortsym} in nine scripts
(\texttt{radiation\_collision/cas}), 128 bindings, each required to reduce to
exactly zero:
the resonance identities, completing the square and the null space
(Sec.~\ref{sec:variational});
the Parseval integrals, the summed integrand \eqref{eq:Deta}, the photon
prefactor, the antiderivatives and end-point evaluation that produce
\eqref{eq:Dclosed}, the tensor transformation from $(p,\chi)$ and the series
\eqref{eq:FHseries} (Sec.~\ref{sec:D});
the common photon, its Jacobian, $P_0$, the resonance-line parametrisation,
the axis limit, the dipole moments, the Bessel-equation identity behind
$B_{k\ell}$ and the coefficients $A_{k\ell},B_{k\ell}$ against Taylor series of
$J_k^2$, $z^2J_k'^2$ (Sec.~\ref{sec:kernel});
the degeneracy integral, the Kapteyn exponent, the $\eta$ integral and
polarisation inequality behind \eqref{eq:tail}, the telescoping tail sum and the
thermal-tail antiderivative (Sec.~\ref{sec:properties});
the threshold expansions, the counter-streaming thresholds and the broadening
substitution (Secs.~\ref{sec:thresholds}--\ref{sec:modified});
the Bessel moments, $G_E$, $G_P$, the angular average, $\sH_E$, $\sH_P$ and the
variational algebra of the closure and the redistribution model
(Secs.~\ref{sec:closure}--\ref{sec:interpretation});
the drag characteristics, trajectory, Jacobian and propagator
(Sec.~\ref{sec:drag}). The antiderivatives in the angular integration were
proposed with SymPy and are verified by fortsym by differentiation.

\paragraph{Numerical checks.} \texttt{numeric\_checks.py} (mpmath, 20--30 digits)
tests the non-algebraic statements independently: harmonic sums by direct
summation, Eq.~\eqref{eq:Dclosed} against the $\eta$ quadrature of the unsummed
harmonic series, the absorption series and $S_M$ against direct integration
along the resonance line, Bessel moments, $G_E$, $G_P$ and the by-parts
identity for $\sH_\alpha$ against two-dimensional quadrature, the tail bound at
sample points (the entry is the largest ratio of tail to bound), the broadened cusp,
the drag trajectory against direct integration of Eq.~(11.2), and the
propagator through particle number and the energy-loss rate. Maximal relative
errors:
\begin{center}
\input{numeric_table}
\end{center}
With $n_e=\nu_0=1$, direct quadrature reproduces
\begin{center}
\input{g_table}
\end{center}
A discretisation should build its Gram matrix from the same quadrature as its
moments and diffusion; analytic $G_E,G_P$ combined with inexact discrete
moments do not conserve to machine precision.

\paragraph{Implementation.} Evaluate $F,G,H$ and $D$ with the series
\eqref{eq:FHseries} near the origin, $G_E,G_P$ with scaled Bessel functions for
large $\mu$, then $\sH_E,\sH_P$, the two moments $b_E,b_P$, and add the two
source functions to the differential operator. A weak-form discretisation with
a consistent mass matrix, discrete gradients and a Gram matrix from the same
diffusion bilinear form preserves the null space at any resolution.

\begin{center}
\small
\begin{tabular}{p{3.4cm}p{4.7cm}p{5.6cm}}
\toprule
Model & Analytic result & Condition\\
\midrule
Original local operator & Convergent all-harmonic kernel; exact $D$ & No finite resummation; global obstruction conditional\\
Finite line width & Analytic smoothing of the cusp & Detailed balance must be kept; kernel not closed\\
Two-mode photons & Rank-two integral part, exact $D$ & Common temperature/drift response; residual is $\|(\mathrm{Id}-P)q_*\|_N^2$\\
Strong redistribution & $O(\kappa^{-1})$ convergence to the closure & Specified added redistribution\\
External thermal bath & Finite $2\times2$ response or differential operator only & Reservoir exchanges energy and momentum\\
Optically thin & Closed characteristics and propagator & Drag only; energy leaves with the photons\\
\bottomrule
\end{tabular}
\end{center}

The central simplification keeps the relativistic electron kinematics and the
full harmonic content of $D$ and restricts only the self-consistent photon
perturbation to its two conservation-related modes. This removes the
mode-by-mode reciprocal normaliser before any photon integral is done, and
leaves only constants expressible through $K_\nu(\mu)$.

%======================================================================
\appendix
\section{Angular integration}\label{app:angular}

For $\uperp>0$ the substitution $x=(u^2\eta-\gamma\upar)/\uperp$ gives
$Q=\uperp^2(1+x^2)/u^2$, $\dd\eta=\uperp\dd x/u^2$ and end points
$x_\pm=(\pm u^2-\gamma\upar)/\uperp$, with $\sqrt{1+x_\pm^2}=u(\gamma\mp\upar)/\uperp$ and
$\arsinh x_\pm=\pm\arsinh u-\arsinh(\upar/\uperp)$. The integrals
$I_{m,p}=\int x^m(1+x^2)^{-p/2}\dd x$ (odd $p$) follow from
\begin{gather}
I_{m,p}=I_{m-2,p-2}-I_{m-2,p},\qquad
I_{0,p}=\frac{x}{(p-2)(1+x^2)^{(p-2)/2}}+\frac{p-3}{p-2}I_{0,p-2},\\
I_{1,p}=-\frac{(1+x^2)^{1-p/2}}{p-2},\qquad I_{0,1}=\arsinh x .
\end{gather}
The boost parameter $\arsinh(\upar/\uperp)$ cancels and the result is
Eq.~\eqref{eq:Dclosed}. The verified form uses $\eta$ directly: for each
integrand $\eta^jP_k/Q^{k/2}$ ($k=5,7$, $j=0,1,2$) the antiderivative is
$(R_{kj}Q^{-(k-2)/2}+C_{kj}L)/D_{kj}$ with polynomials $R_{kj}$, $C_{kj}$, $D_{kj}$
in $(\gamma,\upar,\eta)$, which fortsym checks through the polynomial identity
$R'Q-\tfrac{k-2}2RQ'+CQ^{(k-1)/2}=D\eta^jP_k$. The axis limits follow
continuously.

\section{Absorption series}\label{app:series}

The Bessel power series gives $J_k(z)^2=\sum_\ell(-1)^\ell A_{k\ell}(z/2)^{2k+2\ell}$.
The Bessel equation implies
\begin{equation}
z^2J_k'(z)^2=\frac{z^2}2\frac{\dd^2}{\dd z^2}J_k^2+\frac z2\frac{\dd}{\dd z}J_k^2+(z^2-k^2)J_k^2,
\end{equation}
and equating powers gives $B_{k\ell}$. Expanding $e^{-ax}$, keeping even powers and
using the beta integral gives $M_j$, $T_j$; the factor $x^2$ is kept in $T_j$ rather
than written as $M_j-M_{j+1}$, which avoids one cancellation. For fixed $k,a,b$
the Bessel series and its derivatives converge uniformly on $|z|\le|b|$, and
$|e^{-ax}|\le e^{|a|}$ on $[-1,1]$, so termwise integration is justified. The
coefficient ratio
\begin{equation}
\frac{A_{k,\ell+1}}{A_{k\ell}}=\frac{(2k+2\ell+2)(2k+2\ell+1)}{(\ell+1)(2k+\ell+1)(k+\ell+1)^2}
\end{equation}
tends to zero after multiplication by $(b/2)^2$, giving absolute convergence;
$B_{k\ell}$ has only extra polynomial factors. With $M_j,T_j\le2e^{|a|}$ a
geometric majorant bounds the tail once the ratio is below a chosen number
less than one. For fixed $s>0$, $|\eta|<1$, the least equilibrium energy on
harmonic $k$ obeys
$\exp[-\frac\mu\sigma(ks-|\eta|h_k)]\le\exp[-\mu ks/(1+|\eta|)]$, and the
remaining factors are polynomial in $k$, so the outer absorption sum
converges.

%======================================================================
\begin{thebibliography}{99}
\bibitem{KK96} S.~V. Kasilov and W.~Kernbichler, ``Passive cyclotron current drive efficiency for relativistic toroidal plasmas,'' Phys. Plasmas \textbf{3}, 4115--4127 (1996), doi:10.1063/1.871544. Eqs.~(38)--(43), (46)--(55).
\bibitem{KK96b} W.~Kernbichler and S.~V. Kasilov, ``Passive cyclotron current generation with fish-scale structures,'' Phys. Plasmas \textbf{3}, 4128--4139 (1996), doi:10.1063/1.871545.
\bibitem{Akh62} A.~I. Akhiezer, V.~F. Aleksin, V.~G. Bar'yakhtar and S.~V. Peletminskii, ``Effect of radiation on electron relaxation and plasma electrical conductivity in a strong magnetic field,'' Sov. Phys. JETP \textbf{15}, 386--393 (1962).
\bibitem{Akh63} A.~I. Akhiezer, V.~G. Bar'yakhtar and S.~V. Peletminskii, ``Effect of radiation processes on transport phenomena in a plasma in a strong magnetic field,'' Sov. Phys. JETP \textbf{16}, 1231 (1963).
\bibitem{Pel63} S.~V. Peletminskii, ``Relaxation of photons and plasma electrons in a strong magnetic field,'' Sov. Phys. JETP \textbf{17}, 497--501 (1963).
\bibitem{Bertelli05} N.~Bertelli, M.~Bornatici and F.~Engelmann, ``Electron cyclotron absorption in a high-temperature plasma revisited,'' Phys. Lett. A \textbf{347}, 114--120 (2005), doi:10.1016/j.physleta.2005.06.078.
\bibitem{Ghisellini88} G.~Ghisellini, P.~W. Guilbert and R.~Svensson, ``The synchrotron boiler,'' Astrophys. J. \textbf{334}, L5--L8 (1988), doi:10.1086/185300.
\bibitem{Hirvijoki15} E.~Hirvijoki, J.~Decker, A.~J. Brizard and O.~Embr\'eus, ``Guiding-centre transformation of the radiation--reaction force in a non-uniform magnetic field,'' J. Plasma Phys. \textbf{81}, 475810504 (2015), doi:10.1017/S0022377815000744.
\bibitem{Kennedy21} D.~Kennedy and P.~Helander, ``Coulomb collisions in strongly anisotropic plasmas. I. Cyclotron cooling in electron--ion plasmas,'' J. Plasma Phys. \textbf{87}, 905870107 (2021), doi:10.1017/S0022377820001622.
\bibitem{Ochs24a} I.~E. Ochs, M.~E. Mlodik and N.~J. Fisch, ``Electron tail suppression and effective collisionality due to synchrotron emission and absorption in mildly relativistic plasmas,'' Phys. Plasmas \textbf{31}, 083303 (2024), doi:10.1063/5.0228464.
\bibitem{Ochs24b} I.~E. Ochs, ``Synchrotron-driven instabilities in relativistic plasmas of arbitrary opacity,'' Astrophys. J. \textbf{975}, 30 (2024), doi:10.3847/1538-4357/ad76aa.
\bibitem{Abel08} I.~G. Abel, M.~Barnes, S.~C. Cowley, W.~Dorland and A.~A. Schekochihin, ``Linearized model Fokker--Planck collision operators for gyrokinetic simulations. I. Theory,'' Phys. Plasmas \textbf{15}, 122509 (2008), doi:10.1063/1.3046067.
\bibitem{Sugama09} H.~Sugama, T.-H. Watanabe and M.~Nunami, ``Linearized model collision operators for multiple ion species plasmas and gyrokinetic entropy balance equations,'' Phys. Plasmas \textbf{16}, 112503 (2009), doi:10.1063/1.3257907.
\bibitem{DLMF} NIST Digital Library of Mathematical Functions, \url{https://dlmf.nist.gov/}, Chapters~10 and~12 (\S\S10.2, 10.6, 10.19, 10.29, 10.32, 12.2, 12.5).
\bibitem{fortsym} fortsym, computer algebra for Fortran, \url{https://github.com/lazy-fortran/fortsym}; the checks use the revision recorded in \texttt{radiation\_collision/README.md}.
\end{thebibliography}

\end{document}
